\section{Preliminaries}
\label{sec:prelim}

This section fixes the notation used for the rest of the report. The definitions
are standard, but the height convention is not universal, and every numerical
claim we make depends on which one we adopt. We therefore state it explicitly and
return to the consequences in \Cref{rem:conventions}.

\begin{definition}[Binary search tree]
\label{def:bst}
A binary search tree is a binary tree whose nodes carry keys drawn from a totally
ordered set, such that for every node $v$, every key in the left subtree of $v$ is
smaller than the key at $v$, and every key in the right subtree of $v$ is larger.
\end{definition}

We write $T$ for a tree, $n$ for the number of nodes it holds, and $L(v)$ and
$R(v)$ for the left and right subtrees of a node $v$. The height of a tree is the
number of edges on the longest path from its root down to a leaf. A tree
consisting of one node therefore has height $0$, and we assign the empty tree the
height $-1$, which makes the recurrence in \Cref{eq:heightrec} hold without a
special case:
\begin{equation}
  \label{eq:heightrec}
  h(T) =
  \begin{cases}
    -1, & \text{if } T \text{ is empty},\\
    1 + \max\bigl(h(L(v)),\, h(R(v))\bigr), & \text{if } v \text{ is the root of } T.
  \end{cases}
\end{equation}

Searching is the operation everything else is measured against. A search compares
the target with the key at the root, and either stops or descends into one
subtree, discarding the other. The number of comparisons is therefore at most
$h(T)+1$, which gives \Cref{prop:searchcost}.

\begin{proposition}[Search cost]
\label{prop:searchcost}
A search in a binary search tree of height $h$ performs at most $h+1$ key
comparisons.
\end{proposition}

\Cref{prop:searchcost} is the reason height is the quantity worth controlling. A
structure that keeps the height logarithmic keeps every operation logarithmic,
and one that does not keep the height logarithmic offers no guarantee at all.

\subsection{The balance condition}

The AVL condition is stated at a single node, in terms of the difference between
the heights of its two subtrees.

\begin{definition}[Balance factor]
\label{def:bf}
The balance factor of a node $v$ is
\begin{equation}
  \label{eq:bf}
  \bfac(v) = h\bigl(L(v)\bigr) - h\bigl(R(v)\bigr).
\end{equation}
\end{definition}

A positive balance factor means that the left subtree is the taller one, and a
negative balance factor means that the right subtree is. The AVL condition
restricts the balance factor to three values.

\begin{definition}[AVL tree]
\label{def:avl}
An AVL tree is a binary search tree in which every node $v$ satisfies
$\bfac(v) \in \{-1, 0, +1\}$.
\end{definition}

\Cref{fig:balance} draws the condition at one node. A node whose balance factor
reaches $+2$ or $-2$ violates \Cref{def:avl}, and we call such a node
\emph{unbalanced}. No other value can occur during the operations we describe,
because insertion and deletion each change the height of a subtree by at most one
before the repair runs, so a balance factor moves from a legal value to $\pm 2$
and no further.

\begin{figure}[t]
  \centering
  % The whole AVL invariant, drawn at a single node: two subtrees, two heights,
% and the one-level gap the definition permits between them.
\begin{tikzpicture}[x = 1mm, y = 1mm]

  \node[tnode] (v) at (0, 0) {$v$};

  % subtrees as triangles, the left one deliberately one level taller
  \draw[tedge] (v) -- (-16, -12);
  \draw[tedge] (v) -- ( 16, -12);
  \draw[tedge, fill = rwash] (-16, -12) -- (-27, -40) -- (-5, -40) -- cycle;
  \draw[tedge, fill = white] ( 16, -12) -- (  7, -33) -- (25, -33) -- cycle;

  \node[font = \small] at (-16, -30) {$L$};
  \node[font = \small] at ( 16, -25) {$R$};

  % height brackets, one per side
  \draw[decorate, decoration = {brace, amplitude = 4pt}, draw = rmid]
    (-31, -40) -- (-31, -12)
    node[midway, left = 5pt, tlabel, rotate = 90, anchor = south] {$h_L$};
  \draw[decorate, decoration = {brace, amplitude = 4pt, mirror}, draw = rmid]
    ( 29, -33) -- ( 29, -12)
    node[midway, right = 5pt, tlabel, rotate = 90, anchor = south] {$h_R$};

  \node[tannot] at (0, -50)
    {the definition allows $|h_L - h_R| \le 1$, and nothing looser};

\end{tikzpicture}

  \caption{The AVL condition at one node. The definition permits a gap of one
  level between the two subtree heights, and forbids a gap of two.}
  \label{fig:balance}
\end{figure}

\Cref{tab:notation} collects the symbols.

\begin{table}[t]
  \centering
  \caption{Notation used throughout the report.}
  \label{tab:notation}
  \begin{tabular}{@{}ll@{}}
    \toprule
    Symbol & Meaning \\
    \midrule
    $n$            & number of nodes in the tree \\
    $h$, $h(T)$    & height in edges; a single node has height $0$, the empty tree $-1$ \\
    $L(v)$, $R(v)$ & left and right subtrees of the node $v$ \\
    $\bfac(v)$     & balance factor of $v$, defined in \Cref{eq:bf} \\
    $\Nmin(h)$     & fewest nodes an AVL tree of height $h$ can have \\
    $F(k)$         & Fibonacci number, with $F(0)=0$, $F(1)=1$ \\
    $\varphi$      & golden ratio, $(1+\sqrt{5})/2 \approx 1.618$ \\
    $\lgg{x}$      & logarithm of $x$ to base two \\
    \bottomrule
  \end{tabular}
\end{table}
